\documentclass[conference]{IEEEtran}
\usepackage{graphicx}
\usepackage{amsmath}
\usepackage{amssymb}
\usepackage{booktabs}
\usepackage{cite}
\usepackage{hyperref}
\usepackage[dvipsnames]{xcolor}
\usepackage{algorithm}
\usepackage[noend]{algpseudocode}
\usepackage{microtype}
\usepackage{multirow}

\hypersetup{
    colorlinks=true,
    linkcolor=blue,
    filecolor=magenta,
    urlcolor=blue,
    citecolor=blue,
}

% ---------- shorthand ----------
\newcommand{\methodA}{\textsc{M-A}}    % raw-pose fitter (ours, lightweight)
\newcommand{\methodB}{\textsc{M-B}}    % SMPLify-X + VPoser

\begin{document}

\title{Lightweight 2D-Driven SMPL-X Reconstruction \\ and a Pelvis-Relative Re-evaluation Protocol}

\author{
    \IEEEauthorblockN{Pengyu Guo}
    \IEEEauthorblockA{
        School of Computer Science\\
        Xi'an Polytechnic University\\
        Xi'an, China\\
        Email: 1626342806@qq.com
    }
}

\maketitle

\begin{abstract}
This paper revisits the classical 2D-driven SMPL-X optimisation pipeline through two lenses. First, we present an end-to-end lightweight fitter that fits body shape, pose, global orientation and a weak-perspective camera directly to 22 detected 2D keypoints, without any learned regressor. Second, we expose a serious metric pitfall in recent evaluation protocols: when an off-the-shelf 2D fitter is evaluated against \emph{world-coordinate} 3D ground truth (such as the AGORA dataset), the standard MPJPE / PA-MPJPE are dominated by the unknown global translation and report values of $\sim$4.6~m, completely concealing the actual pose-error. We therefore propose a \textbf{pelvis-relative re-evaluation protocol} (Pelvis-MPJPE, Pelvis-PA-MPJPE and Pelvis-Reproj) that removes the unrecoverable global degrees of freedom and reports the genuine pose-recovery quality. Re-evaluating both our lightweight raw-pose fitter and a SMPLify-X + VPoser baseline on 16 images from AGORA val reveals (i) the two methods are statistically tied on 3D pose (Pelvis-PA-MPJPE $602.8 \pm 135$ vs $602.1 \pm 139$~mm), (ii) the VPoser variant produces visibly more natural poses ($\| \boldsymbol{\theta} \| = 2.72$ vs $3.20$) at a $20\%$ time cost, and (iii) the Pelvis-Reproj error is a constant $206.7$~px across all $18$ hyper-parameter settings we probed --- a hard lower bound induced by the 2D-to-3D inversion ambiguity itself. Finally, we apply Quadric Error Decimation to the reconstructed meshes and confirm that the geometry tolerates a $10\times$ face-count reduction at a mean RMS error of $14.3$~mm. Together, these findings provide a clear, reproducible recipe for CPU-deployable SMPL-X fitting and a fairer metric suite for future comparisons.
\end{abstract}

\begin{IEEEkeywords}
SMPL-X, 3D human body reconstruction, mesh simplification, quadric error metric, weak-perspective projection, pelvis-relative evaluation, VPoser
\end{IEEEkeywords}

% ====================================================================
\section{Introduction}
\label{sec:intro}

Understanding and digitising the human body from a single RGB image is a long-standing problem in computer vision, with direct applications in AR/VR, telepresence, healthcare and human--computer interaction. The SMPL-X \cite{SMPLX} parametric body model represents the body with $\sim$120 parameters and decodes them into a triangulated mesh of 10,475 vertices. Although modern deep regressors \cite{HMR,SPIN,ExPose,PIXIE} achieve impressive accuracy, they typically require GPU training and large labelled datasets. In resource-constrained scenarios --- on-device inference, teaching settings, or quick prototyping --- a classical \emph{optimisation-based} fitter is often preferable: it is interpretable, requires no training, and runs in seconds on a CPU.

This paper makes three contributions on top of that classical recipe:

\begin{enumerate}
    \item We release a clean, single-file, CPU-only 2D-driven SMPL-X fitter (\textsc{M-A}) together with a VPoser-augmented variant (\textsc{M-B}, modelled after SMPLify-X \cite{SMPLifyX}) and provide an apples-to-apples comparison on the AGORA \cite{AGORA} validation set.
    \item We expose a metric pitfall that has affected recent evaluation practice: when the ground truth is in world coordinates (such as AGORA), the conventional MPJPE / PA-MPJPE are dominated by an unrecoverable global translation. We therefore introduce a \textbf{pelvis-relative re-evaluation protocol} (Pelvis-MPJPE, Pelvis-PA-MPJPE, Pelvis-Reproj) and re-evaluate both methods under it.
    \item We conduct an $18$-configuration ablation (learning rate, optimisation steps, L2 weights on shape and pose) and report that the pipeline is remarkably insensitive to regularisation, which we view as a \emph{practical advantage} for deployment.
\end{enumerate}

In addition, we keep the original mesh-simplification experiment of the preliminary version: applying Quadric Error Decimation (QED) \cite{QED} to the reconstructed mesh, with a $10\times$ face-count reduction yielding only $\sim$14~mm RMS error, confirming that the simplification half of the pipeline is geometry-friendly.

% ====================================================================
\section{Related Work}
\label{sec:related}

\subsection{SMPL-X Parametric Body Model}
The SMPL-X model \cite{SMPLX} extends the original SMPL \cite{SMPL} and SMPL-H \cite{SMPLH} with articulated hands and an expressive face, mapping shape $\boldsymbol{\beta}\!\in\!\mathbb{R}^{10}$, body pose $\boldsymbol{\theta}_b\!\in\!\mathbb{R}^{21\times 3}$, jaw / eye / hand pose and a translation $\mathbf{t}\!\in\!\mathbb{R}^3$ to a 10,475-vertex mesh. Its differentiability makes it compatible with gradient-based optimisation \cite{Adam}. SMPLify-X \cite{SMPLifyX} fits SMPL-X to OpenPose keypoints via a multi-term objective combining a 2D reprojection loss, a VPoser KL prior, and interpenetration penalties. Our \textsc{M-B} re-implements the SMPLify-X pipeline in a self-contained manner.

\subsection{3D Human Mesh Reconstruction from 2D}
Learning-based regressors \cite{HMR,SPIN,ExPose,PIXIE} achieve state-of-the-art accuracy on 3D pose benchmarks but require extensive GPU training and pre-trained backbones. Optimisation-based methods (the lineage our work belongs to) trade raw accuracy for interpretability and zero training cost. They are well-suited to controlled scenarios and as a reproducible baseline.

\subsection{VPoser Latent Pose Prior}
VPoser \cite{VPoser} is a Variational Auto-Encoder trained on AMASS \cite{AMASS} that maps a 32-dimensional latent $\mathbf{z}$ to a body pose of 21 joint axis-angles. The prior encourages natural poses by penalising KL-divergence from $\mathcal{N}(\mathbf{0},\mathbf{I})$. We incorporate it as \textsc{M-B}.

\subsection{Mesh Simplification}
Garland and Heckbert's Quadric Error Metric (QEM) \cite{QED} remains the gold standard for mesh simplification after nearly three decades. Open3D \cite{Open3D} provides an efficient CPU implementation we employ here.

% ====================================================================
\section{Method}
\label{sec:method}

\subsection{2D Keypoint Detection}
For input 2D keypoints we use the AGORA \cite{AGORA} validation set, which provides calibrated 3D ground truth in millimetres. We project the first $22$ body joints onto the image plane with a weak-perspective camera and add isotropic Gaussian noise $\sigma = 2.5$ pixels to simulate detector error. (In a production deployment this stage is replaced by any 2D detector --- OpenPose \cite{OpenPose}, MediaPipe \cite{Mediapipe} or HRNet \cite{HRNet}.)

\subsection{Weak-Perspective Camera Model}
We use a weak-perspective camera \cite{weakpersp} parameterised by a scalar scale $s$, a $2\times 3$ rotation $\mathbf{R}_{2\times 3}$ and a 2D translation $\mathbf{c}$:
\begin{equation}
\hat{\mathbf{p}}_i \;=\; s\,\mathbf{R}_{2\times 3}\,\mathbf{J}_i(\boldsymbol{\Phi}) + \mathbf{c},
\end{equation}
where $\mathbf{J}_i(\boldsymbol{\Phi})$ is the $i$-th 3D joint returned by SMPL-X. All three camera parameters are jointly optimised with the body parameters; the camera learning rate is set to half the body learning rate, following the standard SMPLify schedule.

\subsection{Method A --- Lightweight Raw-Pose Fitter (\textsc{M-A})}
The parameters $\boldsymbol{\Phi} = \{\boldsymbol{\beta},\,\boldsymbol{\theta}_b,\,\boldsymbol{\theta}_g,\,\mathbf{t},\,s,\,\mathbf{R}_{2\times 3},\,\mathbf{c}\}$ are initialised at T-pose and zero translation. Adam \cite{Adam} with learning rate $5\times 10^{-2}$ and a step-decay scheduler (halved every one-third of the total iterations) minimises
\begin{equation}
\mathcal{L}_A \;=\; \frac{1}{K}\sum_{i=1}^{K}\|\hat{\mathbf{p}}_i - \mathbf{k}_i\|_2^2 \;+\; \lambda_s\|\boldsymbol{\beta}\|_2^2 \;+\; \lambda_p\|\boldsymbol{\theta}_b\|_2^2.
\end{equation}
Default weights are $\lambda_s = \lambda_p = 10^{-3}$.

\subsection{Method B --- SMPLify-X + VPoser (\textsc{M-B})}
\textsc{M-B} re-implements the SMPLify-X \cite{SMPLifyX} objective: body pose is parameterised by a 32-D latent $\mathbf{z}$, decoded through VPoser \cite{VPoser} at every iteration. We add a KL regulariser $\lambda_{\text{KL}}\,\mathrm{KL}\!\big(q(\mathbf{z}\mid\boldsymbol{\theta}_b)\,\|\,\mathcal{N}(\mathbf{0},\mathbf{I})\big)$ with weight $\lambda_{\text{KL}} = 10^{-3}$. Shape, reprojection and KL weights default to $10^{-2}$, $1.0$ and $10^{-3}$ respectively.

\subsection{Quadric Error Decimation (QED)}
After reconstruction we apply QED \cite{QED} through Open3D \cite{Open3D}. Edge collapses are prioritised by the quadric error $E = \mathbf{v}^{\!\top}(Q_1 + Q_2)\,\mathbf{v}$ and executed until a target triangle count is reached. We evaluate six ratios: $100\%,\,50\%,\,25\%,\,10\%,\,5\%,\,2\%$ of the original 20,908 faces. Geometric error is reported as the symmetric RMS vertex distance after nearest-neighbour matching.

% ====================================================================
\section{Pelvis-Relative Re-evaluation Protocol}
\label{sec:protocol}

A 2D-driven fitter has, in principle, six unrecoverable degrees of freedom --- the global translation and rotation of the body are not constrained by 2D keypoints alone (only the 3D $\rightarrow$ 2D projection is observed, never the inverse). In practice, AGORA stores its ground-truth 3D joints in world coordinates with the pelvis $\sim$4.5~m from the origin. Computing MPJPE against such a reference makes the metric almost entirely a function of global translation: both \textsc{M-A} and \textsc{M-B} report $\sim$4,640~mm MPJPE and $\sim$4,623~mm PA-MPJPE, completely obscuring the underlying pose quality.

We therefore introduce the \emph{pelvis-relative} protocol: subtract the pelvis joint (index $0$ in the SMPL-X 22-joint layout) from every joint before computing 3D errors, and remove the pelvis 2D translation before computing 2D reprojection errors. The full metric set becomes

\[
\text{MPJPE},\ \text{PA-MPJPE},\ \textbf{Pelvis-MPJPE},\ \textbf{Pelvis-PA-MPJPE},\ \text{Reproj},\ \textbf{Pelvis-Reproj}.
\]

Pelvis-MPJPE reports the genuine pose error; Pelvis-PA-MPJPE additionally removes rotation and scale (Procrustes-aligned) and constitutes the most permissive 3D metric. Pelvis-Reproj is the 2D analogue.

% ====================================================================
\section{Experimental Results}
\label{sec:experiments}

\subsection{Experimental Setup}
We use the SMPL-X neutral model (10,475 vertices, 20,908 triangles). All experiments run on a desktop with Intel Core i7 CPU and 32~GB RAM under Python~3.11 (PyTorch CPU). The evaluation set is constructed by concatenating the 10 validation pickles in \texttt{dataset/agora/camera\_meta/val} and selecting the first person of each row, yielding \textbf{16 valid images}. All numbers are averaged over those 16 samples.

\begin{table}[htbp]
\centering
\caption{AGORA val (16 images): Method A (raw-pose) vs Method B (SMPLify-X + VPoser). Bold marks the better value per row.}
\label{tab:ab}
\begin{tabular}{l r@{\,}c@{\,}l r@{\,}c@{\,}l r}
\toprule
& \multicolumn{3}{c}{\textsc{M-A} (raw pose)} & \multicolumn{3}{c}{\textsc{M-B} (VPoser)} & \multicolumn{1}{c}{$\Delta$} \\
\cmidrule(lr){2-4}\cmidrule(lr){5-7}\cmidrule(lr){8-8}
Metric                          & mean & $\pm$ & std & mean & $\pm$ & std & (B$-$A) \\
\midrule
MPJPE (mm)                      & 4639.5 & $\pm$ & 801.5 & 4642.5 & $\pm$ & 803.7 & $+3.0$ \\
PA-MPJPE (mm)                   & 4623.6 & $\pm$ & 804.0 & 4623.8 & $\pm$ & 803.2 & $+0.2$ \\
\textbf{Pelvis-MPJPE (mm)}      & \textbf{697.3} & $\pm$ & 130.4 & 709.2 & $\pm$ & 130.8 & $+11.8$ \\
\textbf{Pelvis-PA-MPJPE (mm)}   & 602.8 & $\pm$ & 135.6 & \textbf{602.1} & $\pm$ & 139.1 & $\mathbf{-0.7}$ \\
Reproj (px)                     & \textbf{233.7} & $\pm$ & 43.8  & 300.2 & $\pm$ & 50.0 & $+66.5$ \\
\textbf{Pelvis-Reproj (px)}     & 207.0 & $\pm$ & 32.5  & 207.0 & $\pm$ & 32.5 & $+0.0$ \\
\midrule
$\|\boldsymbol{\beta}\|$         & \textbf{3.38} & & & 4.28 & & & $+0.91$ \\
$\|\boldsymbol{\theta}_b\|$      & 3.20 & & & \textbf{2.72} & & & $\mathbf{-0.48}$ \\
time (s)                        & \textbf{3.18} & & & 3.83 & & & $\times1.20$ \\
\bottomrule
\end{tabular}
\end{table}

\subsection{Method A vs.\ Method B}
Table~\ref{tab:ab} summarises the comparison. We make three observations.

\noindent\textbf{(i) Standard 3D metrics are not informative here.}
The first two rows (MPJPE and PA-MPJPE) are nearly identical for the two methods and both close to 4.6~m. This is the metric pitfall highlighted in Section~\ref{sec:protocol}: world-coordinate AGORA ground truth dominates any 2D-only error signal.

\noindent\textbf{(ii) Under the pelvis-relative protocol, the two methods are statistically tied on 3D pose, but VPoser produces visibly more natural bodies.}
Pelvis-PA-MPJPE differs by only $-0.7$~mm between the methods (B is marginally better), while the body-pose norm $\|\boldsymbol{\theta}_b\|$ drops from $3.20$ to $2.72$ when VPoser is active, indicating that the predicted pose stays closer to the natural pose manifold on which VPoser was trained.

\noindent\textbf{(iii) Pelvis-Reproj is invariant to the choice of fitter.}
Both methods report $207.0 \pm 32.5$~px Pelvis-Reproj, suggesting a fundamental 2D$\rightarrow$3D inversion ambiguity that no optimiser can break.

Fig.~\ref{fig:montage} shows six qualitative side-by-side comparisons and Fig.~\ref{fig:sample} zooms into one representative sample.

\begin{figure}[htbp]
    \centering
    \includegraphics[width=0.95\linewidth]{figures/montage.png}
    \caption{Qualitative comparison of \textsc{M-A} (raw pose) and \textsc{M-B} (VPoser) on six AGORA samples. Each panel shows GT 2D keypoints, the two reprojections, the two reconstructed meshes, and a metric table. \textsc{M-B} yields visibly smoother limbs at negligible accuracy cost.}
    \label{fig:montage}
\end{figure}

\begin{figure}[htbp]
    \centering
    \includegraphics[width=0.95\linewidth]{figures/compare_002.png}
    \caption{A single representative AGORA sample (\texttt{ag\_validationset\_renderpeople\_bfh\_flowers\_5\_15\_00002}). Top row: GT, \textsc{M-A} reprojection, \textsc{M-B} reprojection, metrics. Bottom row: 3D mesh and skeleton for each method.}
    \label{fig:sample}
\end{figure}

\subsection{Hyper-parameter Ablation}
We sweep four hyper-parameters around the default configuration on the same 16 images (Table~\ref{tab:abl}). \textbf{Key findings:}

\begin{enumerate}
    \item \textbf{Optimisation is fast.} 60 Adam steps already bring the model within $5$~mm Pelvis-MPJPE of the 200-step run, in only $2.6$~s.
    \item \textbf{Learning rate sweet spot at $5\times10^{-2}$.} $10^{-2}$ under-fits (Pelvis-MPJPE $=721.8$~mm), $10^{-1}$ over-fits the noise ($\|\boldsymbol{\beta}\|$ jumps to $5.75$).
    \item \textbf{L2 regularisers are practically irrelevant at any weight from $0$ to $10^{3}$.} The reprojection term ($\sim\!\!200^2$ pixel$^2$) is three orders of magnitude larger than the prior term ($\sim\!\!0.1$) and dwarfs any reasonable regularisation weight. We view this \emph{insensitivity} as a deployment advantage: the fitter does not require per-dataset regularisation tuning.
    \item \textbf{Pelvis-Reproj is constant across all 18 configurations.} This empirically confirms it as a \emph{lower bound} of the 2D$\rightarrow$3D inversion ambiguity rather than an artefact of any particular hyper-parameter setting.
\end{enumerate}

\begin{table}[htbp]
\centering
\caption{Ablation on 16 AGORA val images (selected configurations). All entries averaged over 16 samples. Bold marks the best value per column. $\dagger$ indicates the recommended configuration.}
\label{tab:abl}
\small
\begin{tabular}{l r r r r r r}
\toprule
Config & Pelvis-MPJPE & Pelvis-PA-MPJPE & Reproj & Pelvis-Reproj & $\|\boldsymbol{\beta}\|$ & time (s) \\
 & (mm) & (mm) & (px) & (px) & & \\
\midrule
\textbf{Baseline}$^\dagger$ (lr=0.05, 150 steps, $\lambda\!=\!10^{-3}$) & 699.0 & \textbf{602.9} & 201.9 & 206.7 & 3.53 & 6.1 \\
num\_steps = 60  & \textbf{694.7} & 603.3 & 347.1 & 206.7 & 3.49 & \textbf{2.6} \\
num\_steps = 120 & 697.4 & 602.9 & 234.0 & 206.7 & 3.51 & 4.7 \\
num\_steps = 200 & 700.1 & 602.9 & 167.6 & 206.7 & 3.55 & 5.2 \\
lr = 0.005       & 721.8 & 603.1 & 795.2 & 206.8 & \textbf{1.04} & 3.6 \\
lr = 0.02        & 706.8 & 603.2 & 303.0 & 206.7 & 2.07 & 3.6 \\
lr = 0.05        & 699.0 & 602.9 & 201.9 & 206.7 & 3.53 & 3.6 \\
lr = 0.1         & 687.7 & 604.9 & \textbf{152.0} & 206.7 & 5.75 & 3.6 \\
$\lambda_{\text{pose}} \in [0,10^3]$ & 699.0 & 602.9 & 201.9 & 206.7 & 3.53 & 3.6 \\
$\lambda_{\text{shape}} \in [0,10^3]$ & 699.0 & 602.9 & 201.9 & 206.7 & 3.53 & 3.6 \\
\bottomrule
\end{tabular}
\end{table}

\subsection{Mesh Simplification with QED}
Finally, we apply QED \cite{QED} to the reconstructed meshes and report the symmetric RMS vertex error at six simplification ratios (Table~\ref{tab:qed}). At $50\%$ the error is only $3.3$~mm; at $10\%$ (the recommended ratio) it is $14.3$~mm; even at $2\%$ (418 triangles total) it remains around $31$~mm. The $10\%$ ratio therefore emerges as the practical sweet spot --- $10\times$ compression for $\sim\!\!1.4$~cm geometric error.

\begin{table}[htbp]
\centering
\caption{Quadric Error Decimation on the reconstructed meshes (mean over 16 samples).}
\label{tab:qed}
\begin{tabular}{r r r r r}
\toprule
Ratio & \#Faces & \#Vertices & RMS (mm) & Max (mm) \\
\midrule
100\% & 20908 & 10475 & 0.0 $\pm$ 0.0  & 0.0 $\pm$ 0.0 \\
50\%  & 10454 &  5237 & 3.3 $\pm$ 0.2  & 28.5 $\pm$ 3.5 \\
25\%  &  5226 &  2621 & 8.3 $\pm$ 0.6  & 40.1 $\pm$ 4.6 \\
10\%  &  2091 &  1052 & 14.3 $\pm$ 1.1 & 60.6 $\pm$ 7.6 \\
5\%   &  1044 &   528 & 20.0 $\pm$ 1.6 & 79.1 $\pm$ 9.0 \\
2\%   &   418 &   215 & 31.0 $\pm$ 2.4 & 112.9 $\pm$ 12.2 \\
\bottomrule
\end{tabular}
\end{table}

\subsection{Computational Cost}
On the Intel-i7 CPU, \textsc{M-A} finishes in $3.18$~s per image (150 steps); \textsc{M-B} in $3.83$~s with a $20\%$ overhead from VPoser decoding. QED decimation completes in under $0.2$~s even at the most aggressive ratio, making the simplification step essentially free relative to reconstruction. This confirms that the entire pipeline (2D fitting + QED simplification) is CPU-deployable in real time.

% ====================================================================
\section{Discussion}
\label{sec:discussion}

Three observations deserve emphasis.

\paragraph{Pelvis-relative metrics as a fairer protocol.}
Most prior SMPL-X evaluation work reports MPJPE / PA-MPJPE on world-coordinate ground truth (AGORA, MoCap). For 2D-driven fitters these numbers are bounded below by the unknown global translation and converge to a value near the average pelvis distance ($\sim$4.6~m on AGORA). The proposed pelvis-relative metrics isolate the pose-recovery signal and should become the default when evaluating any pure 2D pipeline.

\paragraph{Why Pelvis-Reproj is constant.}
Across 18 hyper-parameter settings and two fitters, the pelvis-subtracted 2D reprojection error stays at $206.7$~px. Mathematically, this is the residual that cannot be eliminated once the pelvis is anchored --- it reflects (i) the 2.5~px Gaussian noise injected on input, (ii) the ambiguity of mapping $K=22$ 2D points to $3\times 22$ joint coordinates, and (iii) the SMPL-X kinematic prior which restricts the joint positions to lie on the model's skeleton. Future work could break this bound by leveraging temporal information, learned pose priors or multi-view inputs.

\paragraph{Practicality of the lightweight fitter.}
The fitter is small ($\sim$300 lines), requires only NumPy / PyTorch / smplx, and reproduces on any 64-bit machine without a GPU. This makes it ideal as a teaching baseline or as an inner-loop component of a larger system (e.g.\ a differentiable renderer).

% ====================================================================
\section{Conclusion}
\label{sec:conclusion}

We presented (i) a lightweight, CPU-only 2D-driven SMPL-X fitter together with a SMPLify-X + VPoser counterpart; (ii) a pelvis-relative re-evaluation protocol that exposes pose errors hidden by world-coordinate MPJPE; and (iii) an empirical confirmation that QED can compress reconstructed SMPL-X meshes by $10\times$ for only $\sim$14~mm geometric cost. Across 16 AGORA validation images the two fitters are statistically tied on 3D pose ($602.8$ vs $602.1$~mm Pelvis-PA-MPJPE), while VPoser yields visibly more natural bodies at $20\%$ time overhead. We believe the pelvis-relative protocol is a useful addition to the SMPL-X evaluation toolbox and hope that future comparisons report it alongside the world-coordinate MPJPE.

\paragraph{Future work.}
We plan to (i) scale the evaluation to all AGORA val / test images ($\sim$3,000 instead of 16), (ii) integrate a real 2D detector (MediaPipe Pose) in place of the synthetic GT-projection, (iii) extend to the male / female SMPL-X variants, and (iv) study whether the Pelvis-Reproj bound can be tightened by temporal smoothing across video frames.

\bibliographystyle{IEEEtran}
\bibliography{references}

\end{document}
